\section{Antecedentes}

\noindent\textbf{Internacional.} La literatura técnica de los fabricantes de PMMA documenta ventanas de temperatura, tiempos de calentamiento por espesor y enfriamiento bajo restricción \cite{plexiglas_forming,throne2008}. Gilormini \emph{et al.} propusieron ecuaciones constitutivas viscoelásticas no lineales para PMMA cerca de $T_g$ \cite{gilormini2010}; Yan \emph{et al.} caracterizaron el \emph{crazing} bajo esfuerzo y su mitigación por recocido \cite{yan2023}. Comercialmente, Shannon Machines ofrece líneas modulares con hilo resistivo de altura ajustable \cite{shannon_line,shannon_bending} y varios fabricantes asiáticos venden dobladoras automáticas con selección de ángulo \cite{mic_auto}.

\noindent\textbf{Nacional y local.} En Colombia hay tradición de trabajos de grado sobre máquinas dobladoras, casi siempre para metales: doblado de tubo por compresión, rodillos y rotación \cite{uniandes2003}, y dobladoras hidráulicas de tubería redonda desarrolladas en EAFIT \cite{eafit_hidraulica}. Aportan metodología de diseño mecánico y validación dimensional, pero su fenómeno dominante es la plasticidad del metal y no la transición vítrea de un polímero amorfo, por lo que el control térmico queda abierto. EAFIT dispone de taller de mecanizado, laboratorios de electrónica e instrumentación, y del Laboratorio de Mecatrónica, que recibirá el prototipo.
